% Beamer's default font size is 11 points. It is possible to set the
% default font size to any of 8, 9, 10, 11, 12, 14, 17, 20
\documentclass[13pt]{beamer}

\usetheme{Warsaw}
%\usecolortheme{dolphin}
\usecolortheme{seahorse}
%\useinnertheme{circles}
\usefonttheme{professionalfonts}
%\usefonttheme{structurebold}
%\usefonttheme{default}

\setbeamertemplate{navigation symbols}{}
 
\usepackage[english]{babel}
\usepackage[latin1]{inputenc}
\usepackage{times}
\usepackage[T1]{fontenc}
\usepackage{fancybox}

% These things can change:
\newcommand{\baseurl}{http://cell.vtt.fi/latex}
\urlstyle{sf}
\newcommand{\plus}{\makebox[2em]{+}}
\newcommand{\minus}{\makebox[2em]{-}}

% In beamer, defining new commands inside the slides fails.
\newcommand{\warning}[1]{%
\begin{center}
\Large
\shadowbox{
\textbf{#1}}
\end{center}
}

\begin{document}

\logo{\includegraphics[width=1.4cm]{img/vttplain}\quad}
\title{\scalebox{1.5}{\textrm{\LaTeX}} Course 2011}
\subtitle{Part 2: Logical commands}
\author{Arho Virkki}
\institute{\textsc{VTT Technical Research Centre of Finland}}
\date{\scalebox{0.8}{\textcolor{gray}{}}}

\begin{frame}[fragile]
  \titlepage
\end{frame}


\begin{frame}[fragile]\frametitle{Logical commands}
The document is described in abstract level. Logical command define,
for example,
\begin{itemize}
\item sections (\verb|\section{foo}|),
\item environments (\verb|\begin{frame}|),
\item emphasis (\verb|\emph{observe}|),
\item \dots
\end{itemize}
and the interpretation of these commands is controlled by the style sheet
(.sty) and it depends, of course, also on the media (.pdf document, .pdf
slides or .html page).
\medskip

{\footnotesize Observe that LaTeX was first introduces in 1985 as an extension
for TeX (first released in 1978), and HTML5/CSS3, copying the same idea, is called the
``future of the web'' (as of 2011).}
\end{frame}

\begin{frame}[fragile]\frametitle{What this means?}
Setting the style by hand is not a good idea\footnote{Compare
this with doing <div style=``font-style: italic''> \ldots </div> in an html
document instead of using <emph> and defining that environment properly in
the cascading style sheet (css)!}.
\bigskip

Avoid:
{\small
\begin{verbatim}
\textbf{Backround}\\[3mm]
In the future, one of the most striking 
challenges will be...
\end{verbatim}}

Instead, use:
{\small
\begin{verbatim}
\section{Backround}
In the future, one of the most striking 
challenges will be...
\end{verbatim}}
\end{frame}

\begin{frame}[fragile]\frametitle{What this also means?}
\begin{itemize}
\item LaTeX is a sophisticated, and thus also a bit involved environment.
\item You should be writing scratch files in plain text (or with paper and
pencil -- and then feed the papers into scanner to archive them). 
\item Nevertheless, I use LaTeX to write down the potential ideas still
being immature for publication just to make them clean.
\end{itemize}
\end{frame}


\begin{frame}[fragile]\frametitle{Headings}
The levels of headings are
\begin{itemize}
\item[-1] \verb|\part{}|
\item[0] \verb|\chapter{}|
\item[1] \verb|\section{}|
\item[2] \verb|\subsection{}|
\item[3] \verb|\subsubsection{}|
\item[4] \verb|\paragraph{}|
\item[5] \verb|\subparagraph{}|
\end{itemize}
\end{frame}
 
\begin{frame}[fragile]\frametitle{Headings\dots}
\begin{itemize}
\item The levels  \verb|\part{}| and \verb|\chapter{}| are not used with all
document types. These command are used mostly in books.
\item Ending command (analog to ending tag) is not needed\footnote{This
is in contrast to XML, where the elements always consist of the starting and
ending tag and the content between them. For example,\\ \texttt{<tag>} must
be ended with \texttt{<\textbackslash tag>}.}.
\item Command \verb|\setcounter{secnumdepth}{n}| $n\in {0,\ldots,5}$
defines the extent of explicit numbering (in front of the headings).
\end{itemize} 
\end{frame}

\begin{frame}[fragile]\frametitle{Environments}
The most common environments are
\verb|equation|, \verb|displaymath| and \verb|itemize|. Suppose that
\texttt{end1} and \texttt{env2} are two environments. Then
\begin{verbatim}
\begin{env1}
  ...
  \begin{env2}
  ...
  \end{env2}
\end{env1}
\end{verbatim}
is legal. Observe that the environments need to be nested: \texttt{env1} can
not end before \texttt{env2}.
\end{frame}


\begin{frame}[fragile]\frametitle{Environments: Lists}
\begin{minipage}{0.4\textwidth}
\begin{verbatim}
\begin{enumerate}
\item one
\item two
\end{enumerate}
\begin{itemize}
\item[+] pros
\item[--] cons
\end{itemize}
\end{verbatim}
\end{minipage}
yields\hfill
\begin{minipage}{0.4\textwidth}
\begin{enumerate}
\item one 
\item two
\end{enumerate}
\begin{itemize}
\item[+] pros
\item[--] cons
\end{itemize}
\end{minipage}
\end{frame}


\begin{frame}[fragile]\frametitle{Environments: Plain text}
\verb|\begin{verbatim}|\\
\verb|  o o|\\
\verb|   * |\\
\verb|  \_/|\\
\verb|\end{verbatim}|\bigskip

produces
\begin{verbatim}
  o o 
   * 
  \_/
\end{verbatim}
\end{frame}


\begin{frame}[fragile]\frametitle{Environments: Plain text\dots}
Source code listings are convenient to include with
\verb|\verbatiminput{myprog.f90}|.\bigskip

For algorithm, one should use an environment that can emphasize
the keywords.\footnote{This is a good example of an exercise for this course:\\
Find a package for this and describe how it is used.}
\end{frame}

\begin{frame}[fragile]\frametitle{Environments: Text alignment}
\texttt{center}:
\begin{center}
$\sum$
\end{center}\bigskip

\texttt{flushright}:
\begin{flushright}
$\sum$
\end{flushright}\bigskip

\texttt{flushleft}:
\begin{flushleft}
$\sum$
\end{flushleft}
\end{frame}


\begin{frame}[fragile]\frametitle{Environments: Text alignment\dots}
\begin{verbatim}
\begin{center}
KYNTTIL�T SYTTYV�T VARHAIN\bigskip

Kiert�v�t unettavat auringonnousut\\
Kaikkialle harsona niin\\
Aattoni v�hiin k�ynyt\\
Kun kev�t uutena nousee\\
\dots
\end{center}
\begin{flushright}
\emph{-- Kuusumun profeetta}
\end{flushright}
\end{verbatim}
\end{frame}


\begin{frame}[fragile]\frametitle{Environments: Text alignment\dots}
\begin{center}
KYNTTIL�T SYTTYV�T VARHAIN\bigskip

Kiert�v�t unettavat auringonnousut\\
Kaikkialle harsona niin\\
Aattoni v�hiin k�ynyt\\
Kun kev�t uutena nousee\\
\dots
\end{center}
\begin{flushright}
\emph{-- Kuusumun profeetta}
\end{flushright}
\end{frame}


\begin{frame}[fragile]\frametitle{Environments: Text alignment\dots}
The 'quote' environment is rather basic -- but it works.
\begin{verbatim}
\begin{quote}
He has a profound respect for old age. 
Especially when it's bottled.
--- Gene Fowler 
\end{quote}
\end{verbatim}
\begin{quote}
He has a profound respect for old age. Especially when it's bottled.
--- Gene Fowler
\end{quote}\medskip

The command \texttt{quotation} works better for longer quotations.
\end{frame}


\begin{frame}[fragile]\frametitle{Environments: Table}
\begin{verbatim}
\begin{tabular}{|l|c|c|}
\hline
Name & J. Foo & G. Bar \\
\hline
\hline
A-score & 1 & 3 \\
B-score & -2 & 0 \\
\hline
\end{verbatim}
\end{frame}


\begin{frame}[fragile]\frametitle{Environments: Table\dots}
\begin{center}
\begin{tabular}{|l|c|c|}
\hline
Name & J. Foo & G. Bar \\
\hline
\hline
A-score & 1 & 3 \\
B-score & -2 & 0 \\
\hline
\end{tabular}
\end{center}\bigskip

Somehow this looks shomehow crowded\dots Let's change the title
into {\small
\begin{verbatim}
\rule[-4pt]{0pt}{16pt}Nimi & J. Foo & G. Bar \\
\end{verbatim}
} 
\end{frame}


\begin{frame}[fragile]\frametitle{Environments: tabular\dots}
\begin{center}
\begin{tabular}{|l|c|c|}
\hline
\rule[-4pt]{0pt}{16pt}Name & J. Foo & G. Bar \\
\hline
\hline
A-score & 1 & 3 \\
B-score & -2 & 0 \\
\hline
\end{tabular}
\end{center}\bigskip
\begin{flushleft}
We extended the borders of the table with a line of with zero.\medskip

The \LaTeX{} tables are, to be honest, rather clumsy. This is not a problem
if we can generate them automatically (e.g. from the R language using the 
xtable package).
\end{flushleft}
\end{frame}


\begin{frame}[fragile]\frametitle{Environments: Custom definitions}
Let us define an environment than adds automatically a box at the end of a
proof, and writes the title 'New Proof.' in the beginning: \bigskip
\begin{verbatim}
\newenvironment{newproof}
{\makebox[2cm][l]{\textbf{New Proof.\ }}}
{\hfill $\Box$}
\end{verbatim}
\end{frame}

\begin{frame}[fragile]\frametitle{Environments\dots}
\begin{verbatim}
\begin{newproof}
Suppose that $\epsilon > 0$ is already 
chosen,\dots
\end{newproof}
\end{verbatim}\bigskip

\newenvironment{newproof}
{\makebox[2cm][l]{\textbf{New Proof.\ }}}
{\hfill $\Box$}
\begin{newproof}
Suppose that $\epsilon > 0$ is already chosen, \dots
\end{newproof}
\end{frame}


\begin{frame}[fragile]\frametitle{Environments\dots}
The general form of the environment definition:
\begin{verbatim}
\newenvironment{name}[args]{begdef}{enddef}  
\end{verbatim}
or
\begin{verbatim}
\renewenvironment{name}[args]{begdef}{enddef}
\end{verbatim}
if we wish to re-define an existing environment.
\end{frame}


\begin{frame}[fragile]\frametitle{Environments\dots}
These exists a command for stating lemmas:
\begin{verbatim}
\newtheorem{command name}{theorem name}[counter]
\end{verbatim}
For example, 
\begin{verbatim}
\newtheorem{newlemma}{Lemma}[page]

\begin{newlemma}[The Professors' Remainder 
Theorem]
If we use half of the time left for each 
exercise, we can deal with an infine number 
of exercises. 
\end{newlemma}
\end{verbatim}
\end{frame}

\begin{frame}[fragile]\frametitle{Environments\dots}
\newtheorem{newlemma}{Lemma}[page]
\begin{newlemma}[The Professors' Remainder 
Theorem] 
If we use half of the time left for each 
exercise, we can deal with an infine number 
of exercises.
\end{newlemma}\bigskip

Let us further inspect the command:
\begin{verbatim}
\newtheorem{newlemma}{Lemma}[page]
\end{verbatim}
where the numbering was set to follow the
\texttt{page} counter. Thus, (if this was an article and not a slideshow)
all formulas are numbered as
\begin{center} 
(page,lemma) 
\end{center}
\end{frame}


\begin{frame}[fragile]\frametitle{Environments\dots}
\begin{verbatim}
\newtheorem{mynote}{Note!}
\begin{mynote}
Huom\dots
\end{mynote}
\end{verbatim}\bigskip

\newtheorem{mynote}{Note!}
\begin{mynote}
Notes lose their meaning if there are too many of them, or the notes
are self-evident.
\end{mynote}
\end{frame}


\begin{frame}[fragile]\frametitle{Environments\dots}

\begin{mynote}
\LaTeX{} has lots of counters that can be cited with the command
\verb|\the<counter_name>|:
\begin{verbatim}
We are on slide \thepage.
\end{verbatim}
We are on slide \thepage.
\end{mynote}
\end{frame}


\begin{frame}[fragile]\frametitle{Environments: Summary}
Remember the commands
\begin{itemize}
\item \verb|\newcommand| and
\item \verb|\newtheorem|.
\end{itemize}
Details can be found from books or on-line help systems.\medskip

You can inspect e.g. 
\url{http://en.wikibooks.org/wiki/LaTeX}, \\
or by typing the command name into Google --- which seems to be a reasonable
default procedure for anything nowadays.
\end{frame}


\begin{frame}[fragile]\frametitle{Commands}
If we need a simple command, e.g. 
\verb|\warning{text}|, one can define
\begin{verbatim}
\newcommand{\warning}[1]{%
\begin{center}
\Large
\shadowbox{
\textbf{#1}}
\end{center}
}
\end{verbatim}
\end{frame}


\begin{frame}[fragile]\frametitle{Commands\dots}
Now, the command
\begin{verbatim}
\warning{orthogonal vectors $\neq$ 
orthonormal vectors}
\end{verbatim}\medskip
produces
\warning{orthogonal vectors $\neq$ orthonormal vectors}
\end{frame}


\begin{frame}[fragile]\frametitle{Observations}
\begin{itemize}
\item \verb|\shadowbox{}| is from the package 
\texttt{fancybox}. There are tons of similar useful and less useful packages
available at the Comprehensive TeX Archive Network (CTAN).
\item To define new high level logical commands, one of course needs to know
something about the low level \LaTeX{} ``programming''.   
\end{itemize}
\end{frame}


\end{document}
